\begin{scanpage}{0070}{53}
Substituting here the explicit expressions for the functions $f$ and $g$, we arrive at the final expressions for the fringing field
\[\left.\begin{aligned}E_\phi=-H_\theta&=-\frac{iaH_{0x}}2\left[\frac{J_1(ka\sin\theta)}{\sin\frac\theta2}-i\frac{J_2(ka\sin\theta)}{\cos\frac\theta2}\right]\frac{e^{ikR}}R\cos\phi,\\E_\theta=H_\phi&=-\frac{iaH_{0x}}2\left[\frac{J_1(ka\sin\theta)}{\sin\frac\theta2}+i\frac{J_2(ka\sin\theta)}{\cos\frac\theta2}\right]\frac{e^{ikR}}R\sin\phi.\end{aligned}\right\}\tag{9.03}\]
These equations are valid in the right half-space $\left(0<\theta<\frac\pi2\right)$. In the left half-space $\left(\frac\pi2<\theta<\pi\right)$, the fringing field is easily found by assuming that its electric field is an even function, and its magnetic field an odd function of the $z$ coordinate:
\[\left.\begin{aligned}E_\phi(z)&=E_\phi(-z),\\H_\phi(z)&=-H_\phi(-z).\end{aligned}\right\}\tag{9.04}\]
Consequently, in the region $z<0$ (that is, when $\frac\pi2<\theta<\pi$)
\[\left.\begin{aligned}E_\phi=-H_\theta&=-\frac{iaH_{0x}}2\left[\frac{J_1(ka\sin\theta)}{\cos\frac\theta2}-i\frac{J_2(ka\sin\theta)}{\sin\frac\theta2}\right]\frac{e^{ikR}}R\cos\phi,\\E_\theta=H_\phi&=\frac{iaH_{0x}}2\left[\frac{J_1(ka\sin\theta)}{\cos\frac\theta2}+i\frac{J_2(ka\sin\theta)}{\sin\frac\theta2}\right]\frac{e^{ikR}}R\sin\phi.\end{aligned}\right\}\tag{9.05}\]
Assuming that in Equations (9.03) and (9.05) $\theta=0$ and $\theta=\pi$, respectively, we obtain
\[E_y=-\frac{ika^2}2H_{0x}\frac{e^{ikR}}R,\quad E_x=0,\tag{9.06}\]
which is equivalent to the physical optics approach [see Equation (7.20)].
\end{scanpage}
